% TODO make hw stylesheet
\documentclass[10pt]{article}
\usepackage[letterpaper,top=1in,left=1in,right=1in]{geometry}
\usepackage[dvipsnames,svgnames]{xcolor}
\usepackage[shortlabels]{enumitem}
\usepackage[framemethod=TikZ]{mdframed}
\usepackage{amsmath,amssymb,amsthm}
\usepackage[colorlinks]{hyperref}
\usepackage{multicol}
\usepackage{tabularx}
\usepackage{bbm}

\hypersetup{
    colorlinks=true,
    linkcolor=blue,
    filecolor=magenta,
    urlcolor=magenta,
    pdftitle={MAT 344 HW}
}

\usepackage{mathtools}
\usepackage{thmtools}
\usepackage[textsize=footnotesize]{todonotes}
\usepackage{titling}
\usepackage{cancel}

\reversemarginpar
\mdfdefinestyle{mdbluebox}{roundcorner=10pt,innerbottommargin=9pt,
linecolor=blue,backgroundcolor=TealBlue!5,}
\declaretheoremstyle[headfont=\sffamily\bfseries\color{MidnightBlue},
mdframed={style=mdbluebox},]{thmbluebox}

\mdfdefinestyle{mdbloobox}{frametitlefont=\bfseries,innerbottommargin=8pt,
nobreak=true,backgroundcolor=TealBlue!5,linecolor=blue,}
\declaretheoremstyle[headfont=\bfseries\color{MidnightBlue},
mdframed={style=mdbloobox},headpunct={\\[3pt]},postheadspace=0pt,]{thmbloobox}

\mdfdefinestyle{mdredbox}{frametitlefont=\bfseries,innerbottommargin=8pt,
nobreak=true,backgroundcolor=Salmon!5,linecolor=RawSienna,}
\declaretheoremstyle[headfont=\bfseries\color{RawSienna},
mdframed={style=mdredbox},headpunct={\\[3pt]},postheadspace=0pt,]{thmredbox}

\mdfdefinestyle{mdgreenbox}{linecolor=ForestGreen,backgroundcolor=ForestGreen!5,
linewidth=2pt,rightline=false,leftline=true,topline=false,bottomline=false,}
\declaretheoremstyle[headfont=\bfseries\sffamily\color{ForestGreen!70!black},
mdframed={style=mdgreenbox},headpunct={ --- },]{thmgreenbox}

\mdfdefinestyle{mdorangebox}{linecolor=RedOrange,backgroundcolor=RedOrange!5,
linewidth=2pt,rightline=false,leftline=true,topline=false,bottomline=false,}
\declaretheoremstyle[headfont=\bfseries\sffamily\color{RedOrange!70!black},
mdframed={style=mdorangebox},headpunct={ --- },]{thmorangebox}

\mdfdefinestyle{mdblackbox}{linecolor=black,backgroundcolor=RedViolet!5!gray!5,
linewidth=3pt,rightline=false,leftline=true,topline=false,bottomline=false,}
\declaretheoremstyle[mdframed={style=mdblackbox}]{thmblackbox}

\declaretheoremstyle[spaceabove=6pt,spacebelow=6pt]{boldhead}
\declaretheoremstyle[spaceabove=6pt,spacebelow=6pt,headfont=\normalfont\itshape]{italichead}

\declaretheorem[style=boldhead,name=Problem]{problem}
\declaretheorem[style=boldhead,name=Problem,numbered=no]{problem*}
\declaretheorem[style=boldhead,name=Setup]{setup} % for config geo
\declaretheorem[style=boldhead,name=Example,sibling=problem]{example}
\declaretheorem[style=boldhead,name=Theorem,sibling=problem]{theorem}
\declaretheorem[style=boldhead,name=Corollary,sibling=theorem]{corollary}
\declaretheorem[style=boldhead,name=Proposition,sibling=theorem]{prop}
\declaretheorem[style=boldhead,name=Lemma,numberwithin=problem]{lemma}
\declaretheorem[style=boldhead,name=Theorem,numbered=no]{theorem*}
\declaretheorem[style=boldhead,name=Lemma,numbered=no]{lemma*}
\declaretheorem[style=boldhead,name=Claim,numberwithin=problem]{claim}
\declaretheorem[style=boldhead,name=Claim,numbered=no]{claim*}
\declaretheorem[style=boldhead,name=Remark,numberwithin=problem]{remark}
\declaretheorem[style=boldhead,name=Remark,numbered=no]{remark*}
\declaretheorem[style=boldhead,name=Definition,sibling=theorem]{definition}
\declaretheorem[style=boldhead,name=Definition,numbered=no]{definition*}

\providecommand{\ol}{\overline}
\providecommand{\ceil}[1]{\left\lceil #1 \right\rceil}
\providecommand{\floor}[1]{\left\lfloor #1 \right\rfloor}
\newcommand{\dd}{\mathrm{d}}
\providecommand{\ul}{\underline}
\providecommand{\eps}{\varepsilon}
\providecommand{\half}{\frac{1}{2}}
\providecommand{\CC}{\mathbb C}
\providecommand{\FF}{\mathbb F}
\providecommand{\II}{\mathbb I}
\providecommand{\NN}{\mathbb N}
\providecommand{\QQ}{\mathbb Q}
\providecommand{\RR}{\mathbb R}
\providecommand{\ZZ}{\mathbb Z}
\providecommand{\ind}{\mathbbm 1}
\providecommand{\dg}{^\circ}
\providecommand{\ii}{\item}
\providecommand{\setdiff}{\smallsetminus}
\providecommand{\alert}[1]{{\sffamily\textbf{\textcolor{blue}{#1}}}}
\providecommand{\inv}{^{-1}}
\providecommand{\mb}[1]{\mathbf{#1}}
\newcommand{\what}{\widehat}

\newenvironment{soln}{\begin{proof}[Solution]}{\end{proof}}

\providecommand{\pow}{\operatorname{Pow}}
\providecommand{\vocab}[1]{{\textbf{\textcolor{ForestGreen}{#1}}}}
\providecommand{\lcm}{\operatorname{lcm}}
\providecommand{\abs}[1]{\left\lvert #1\right\rvert}
\providecommand{\smabs}[1]{\lvert #1\rvert}
\providecommand{\innerprod}[1]{\langle\!\langle #1\rangle\!\rangle}
\providecommand{\norm}[1]{\left\| #1\right\|}
\providecommand{\smnorm}[1]{\| #1\|}
\providecommand{\gr}{\operatorname{gr}}

\usepackage{mathrsfs}

% place title at top right
\setlength{\droptitle}{-8ex}
\pretitle{\begin{flushright}\Large\bfseries}
\posttitle{\par\end{flushright}}
\preauthor{\begin{flushright}}
\postauthor{\end{flushright}}
\predate{\begin{flushright}}
\postdate{\end{flushright}}

\title{MAT 344 HW03}
\author{\large Aditya Pahuja}
\date{\large Due 02/17}
\begin{document}
\maketitle

\begin{problem}
    Let $I = [0,1]$ and consider a smooth map
    $H\colon I\times I\to\Omega\subseteq\CC$
    where $\Omega$ is a good domain.
    We consider $H$ as a homotopy of curves
    $\gamma_t\colon I\to\Omega$
    where $\gamma_t(s) = H(t,s)$ for $s\in I$.
    We fix $z_0\in\Omega$ and assume that $\gamma_t(0)=\gamma_t(1)=z_0$
    so $\gamma_t$ is a family of loops in $\Gamma$ based at $z_0$.
    \begin{enumerate}[(i)]
	\ii Let $\eps = \min_{(t,s)\in I\times I}
	\{\operatorname{distance}(H(t,s),\partial\Omega)\}$.
	Why is $\eps > 0$?
	\ii Prove there exists $\delta > 0$ such that
	if two points $p,q\in I\times I$ have $\abs{p-q} < \delta$
	then $\abs{H(p) - H(q)} < \eps/4$.
	\ii Let $\omega$ be a $1$-form with $d\omega = 0$.
	Show that the line integrals $\int_{\gamma_t}\omega$
	are independent of $t\in[0,1]$.
    \end{enumerate}
\end{problem}

\begin{soln}
    (i) The function $d\colon (I\times I)\times\partial\Omega\to[0,\infty)$
    such that $d((t,s),p) = \abs{H(t,s)-p}$ is continuous.
    Since $I\times I$ and $\Omega$ are compact,
    their product is compact, so $d$ achieves a minimum value
    which is equal to $\eps$ by definition.
    Assuming that $H$ maps $I\times I$ into the interior of $\Omega$ so that
    $H(t,s)$ doesn't lie on $\partial\Omega$ for any $(t,s)$,
    $d((t,s),p) > 0$ everywhere, so its minimum must be greater than zero.

    (ii) This follows from the fact that $H\colon I\times I\to\Omega$
    is continuous with compact domain, hence uniformly continuous.

    (iii) Let $n$ be a positive integer such that $1/n < \delta/2$.
    For each $j,k\in\{1,\dots,n\}$, the squares
    $R_{j,k} = [(j-1)/n,j/n]\times[(k-1)/n,k/n]$
    are each contained in an open ball
    of radius $\eps/4$ that is contained in the interior of $\Omega$,
    which are in turn contained in squares $S_{j,k}$ of radius $\eps/2$
    with side lengths parallel to the axes
    that are still contained in the interior of $\Omega$.
    That is, $S_{j,k}$ contains $H(R_{j,k})$.

    For a fixed $j$ and $k\in\{0,1,\dots,n\}$, let
    $a_k = H((j-1)/n, k/n)$ and $b_k = H(j/n,k/n)$.
    Define $\gamma^{j,k}$ to be the restriction of
    $\gamma_{j/n}$ to $[(k-1)/n,k/n]$
    and let $\gamma^j$ be the path $\gamma_{j/n}$, for brevity.
    Let $\eta_k\colon I\to\CC$ be the line segment going from $a_k$ to $b_k$
    for $k\in\{0,1,\dots,n\}$ and $\tilde\eta_k$ be the segment
    going in the other direction, from $b_k$ to $a_k$.
    We argued in Tuesday's class that
    line integrals inside the $S_{j,k}$ depend only on
    the endpoints of the path being integrated over, so we can compute
    \begin{align*}
	\int_{\gamma^{j-1}}\omega
	&= \sum_{k=1}^n \int_{\gamma^{j-1,k}}\omega
	= \sum_{k=1}^n \left(\int_{\eta_{k-1}}\omega
	+ \int_{\gamma^{j,k}}\omega + \int_{\tilde\eta_k}\omega\right)\\
	&= \sum_{k=1}^n\left(\int_{\gamma^{j,k}}\omega
	+ \int_{\tilde\eta_k}\omega + \int_{\eta_k}\omega\right)
	= \sum_{k=1}^n\left(\int_{\gamma^{j,k}}\omega\right)
	= \int_{\gamma^j}\omega
    \end{align*}
    hence $\int_{\gamma^0}\omega = \int_{\gamma^1}\omega$
    as desired.
\end{soln}

\pagebreak
\begin{problem}
    Consider the holomorphic $1$-form
    \begin{equation*}
	\omega = \frac 1{z^2 - 1}\dd z.
    \end{equation*}
    Is this the derivative of a holomorphic function in
    \begin{enumerate}[(i)]
	\ii $\CC - \{-1,1\}$?
	\ii $\CC - [-1,1]$?
	\ii $\CC - (-\infty,1]-[1,\infty)$?
    \end{enumerate}
    Explain your answers.
\end{problem}

\begin{soln}
    (i) \framebox{No}. By Cauchy's integral formula
    applied to the function $f(z) = \frac 1{z+1}$,
    \begin{equation*}
	\int_{B_{1/2}(1)} \omega = 2\pi if(1) = \pi i\ne 0,
    \end{equation*}
    but this integral should be zero if $f$ has an antiderivative.
    \\

    The answer to both (ii) and (iii) is \framebox{yes},
    which are special cases of the following lemma:
    \begin{lemma}
        Let $\omega = f\dd z$ be a holomorphic $1$-form defined on
	$U\subseteq\CC$ an open subset. Let $z_0\in U$ be a fixed point.
	Suppose that for every loop $\gamma\colon I\to U$
	with $\gamma(0) = \gamma(1) = z_0$,
	\begin{equation*}
	    \int_\gamma\omega = 0,
	\end{equation*}
	i.e. line integrals of $\omega$ depend only on the endpoints
	of the path being integrated over.
	Then $f$ has an antiderivative defined by
	\begin{equation*}
	    F(z) = \int_{\gamma_z}\omega
	\end{equation*}
	where $\gamma_z$ is an arbitrary path from $z_0$ to $z$.
    \end{lemma}

    \begin{proof}
        Let $w_0\in U$ and pick $\eps>0$ such that $f$
	is given by a convergent power series
	$f(z) = \sum_{n=0}^\infty a_n(z-w_0)^n$
	in a neighborhood of $B_\eps(w_0)\subseteq U$,
	and let $\gamma_{w_0,w}\colon I\to U$
	defined by $\gamma_{w_0,w}(t) = w_0 + (w-w_0)t$ be
	the line segment from $w_0$ to $w$ for any $w\in B_\eps(w_0)$. Then,
	\begin{align*}
	    \frac{F(w) - F(w_0)}{w-w_0}
	    &= \frac 1{w-w_0}\int_{\gamma_{w_0,w}}
	    \sum_{n=0}^\infty a_n(z-w_0)^n\dd z\\
	    &= \sum_{n=0}^\infty \int_{\gamma_{w_0,w}}
	    a_n(z-w_0)^n\cdot \frac 1{w-w_0}\dd z\\
	    &= \sum_{n=0}^\infty \int_0^1
	    a_n((w-w_0)t)^n\cdot \frac 1{w-w_0}\cdot(w-w_0)\dd z\\
	    &= \sum_{n-0}^\infty \frac{a_n}{n+1}(w-w_0)^n,
	\end{align*}
	where the sum and integral commute because
	power series converge locally uniformly.
	Then, we can take the limit as $w\to w_0$
	to obtain $a_0 = f(w_0)$ (swapping sum and limit
	by local uniform convergence),
	so $F$ is indeed holomorphic at every $w_0\in U$
	with derivative $f$.
    \end{proof}

    Now we show that the given $1$-form $\omega$
    satisfies the hypotheses of the above lemma
    in the case of both (ii) and (iii).
    \\
    
    (ii) Let $\gamma_1$ be the square of side length $1$ centered at $1$,
    $\gamma_{-1}$ the square of side length $1$ centered at $-1$,
    and $\gamma$ the $1\times 2$ rectangle bounding the two squares.
    By Cauchy's integral formula applied to $\frac 1{z+1}$ and to
    $\frac 1{z-1}$ over $\gamma_1$ and $\gamma_{-1}$, respectively,
    we can compute that
    \begin{equation*}
	\int_\gamma\omega = \int_{\gamma_1}\omega
	+ \int_{\gamma_{-1}}\omega
	= 2\pi i\cdot\left(\frac 1{1+1} + \frac 1{-1-1}\right) = 0.
    \end{equation*}
    Let $z_0 = 3/2$ be a point on $\gamma$.
    The fundamental group $\pi_1(\CC-[-1,1],z_0)$ is
    the infinite cyclic group generated by the homotopy class of $\gamma$,
    so any loop based at $z_0$ is homotopic to a multiple of $\gamma$
    through a homotopy fixing $z_0$; thus the integral of $\omega$
    over any loop based at $z_0$ is zero.

    (iii) Any loop based at $z_0 = 0$ is nullhomotopic
    because $U = \CC - (-\infty,-1] - [1,\infty)$ is simply connected;
    in fact, $U$ is contractible via the homotopy
    $H\colon U\times I\to U$ given by $H(p,t) = tp$.
\end{soln}

\begin{problem}
    Let $\Omega\subseteq\CC$ be a good domain
    and $z_1,\dots,z_k\in\operatorname{Int}\Omega$.
    Suppose that $f\colon\Omega - \{z_1,\dots,z_k\}\to\CC$
    is holomorphic and satisfies
    \begin{equation*}
	\lim_{z\to z_j} (z-z_j) f(z) = 0
    \end{equation*}
    for $j = 1,\dots,k$. Show that
    \begin{equation*}
	f(z) = \frac 1{2\pi i}\int_{\partial\Omega}
	\frac{f(\zeta)}{\zeta-z}\dd\zeta
    \end{equation*}
    for $z\in\Omega-\{z_1,\dots,z_k\}$.
\end{problem}

\begin{soln}
    Given $j\in\{1,\dots,k\}$, let $\eps>0$ be such that
    $B_\eps(z_j)$ is contained in the interior of $\Omega$
    and does not contain any other $z_i$.
    On this open ball, define $g_j(z) = (z-z_j)^2f(z)$
    for $z\ne z_j$ and $g_j(z_j) = 0$.
    This is of course holomorphic on $B_\eps(z_j) - \{z_j\}$ and
    \begin{equation*}
	\lim_{z\to z_j} \frac{g_j(z) - g_j(z_j)}{z - z_j}
	= \lim_{z\to z_j} (z-z_j)f(z) = 0,
    \end{equation*}
    so $g_j$ is holomorphic on $B_\eps(z_j)$.
    Thus $g_j$ has a power series representation
    $g_j(z) = \sum_{n=0}^\infty a_n(z-z_j)^n$.
    Since $g_j(z_j) = 0$, $a_0 = 0$,
    and $g_j'(z_j) = 0$, so $a_1 = 0$. Thus
    \begin{equation*}
	\frac{g_j(z)}{(z-z_j)^2}
	= \sum_{n=0}^\infty a_{n+2}(z-z_j)^n
    \end{equation*}
    equals $f(z)$ on $B_\eps(z_j) - \{z_j\}$,
    meaning $f$ extends to a holomorphic function $g_j(z)$
    on the entirety of $B_\eps(z_j)$.
    We have therefore shown that there exists a holomorphic function
    $\tilde f\colon\Omega\to\CC$ such that
    $\tilde f(z) = f(z)$ for $z\in\Omega-\{z_1,\dots,z_k\}$,
    so, by Cauchy's integral formula,
    \begin{equation*}
	\frac 1{2\pi i} \int_{\partial\Omega}
	\frac{f(z)}{\zeta-z}\dd\zeta
	= \frac 1{2\pi i} \int_{\partial\Omega}
	\frac{\tilde f(z)}{\zeta-z}\dd\zeta
	= \tilde f(z) = f(z)
    \end{equation*}
    for all $z\in\Omega-\{z_1,\dots,z_k\}$.
\end{soln}










\end{document}
